
\section{Protectionism}



\begin{frame}{Governance of the Open Economy}
    \mytext{
        With trade, the tools of the state become less effective. 
        
        The state may want to influence the outcome. 

        \begin{itemize}
            \item controls at the border
            \begin{itemize}
                \item migration (labour)
                \item capital (?)
                \item goods and services
            \end{itemize}
            \item tariffs, quotas
            \item regulation, other barriers
            \item subsidies
            \item tax competition
            \item monetary policy
            \item contracts, diplomacy
        \end{itemize}
    }
    \mycomment{
        In an open economy, all government actions have an impact not only on the economy, but also abroad; or it may become less effective or ineffective in an open economy. This depends on the (relative) size of the economy (or state). 

        The state can direct policy directly to the cross-border flows, such as tariffs. But also policy that is directed at the "inside" such as taxes, regulations or subsidies, will effect cross-border flows.
    
        % Policies to influence market outcomes are typically limited in scope. Sometimes, they may not be successful at all, have detrimental effects, or side effects.
    }
\end{frame}


\begin{frame}{Tariffs}
    \mycomment{
    Governments have different means to control imports and exports. A popular version is to raise tariffs. Without, all domestic demand is served by domestic supply.
    }
    \showmytikz[1]{../shared/tikz/tikz-tariff}
    
\end{frame}

\begin{frame}{Tariffs}
    \mycomment{
    If we open the economy to trade, the domestic market is facing the world market price for the good. If this is lower, importers undercut a part of the domestic suppliers. The result is import. 
    }
    \showmytikz[2]{../shared/tikz/tikz-tariff}
\end{frame}

\begin{frame}{Tariffs}
    \mycomment{
    Openness to imports have pivoted the supply curve. The result is an increase in consumer surplus to the detriment of local producers. Now these producers are likely to lobby for a tariff to protect their industry.
    }
    \showmytikz[1,2,3]{../shared/tikz/tikz-tariff}
\end{frame}

\begin{frame}{Tariffs}
    \mycomment{
    A tariff will rise the price of the imported goods. The effect will be the same as an increase of the world market price. This will reduce imports. 

    Not that there can be a tariff that is prohibitively high that would raise the supply curve above $E_1$.
    }
    \showmytikz[1,2,3,4]{../shared/tikz/tikz-tariff}
\end{frame}

\begin{frame}{Tariffs}
    \mycomment{
    In the new situation, imports have gone down, as has consumer surplus. 
    }
    \showmytikz[1,2,3,4,5]{../shared/tikz/tikz-tariff}
\end{frame}

\begin{frame}{Tariffs}
    \mycomment{
    The government has gained an income from the taxes. Note that there will be a certain tariff that maximises government revenue; this will be above zero but below the prohibitive tariff. 

    Today, only a fraction of government revenue stems from tariffs.
    }
    \showmytikz[1,2,3,4,5,6]{../shared/tikz/tikz-tariff}
\end{frame}



\clearpage

\begin{frame}{Import Quotas}
    \showmytikz[1,2]{../shared/tikz/tikz-quota.tex}
    \mycomment{
        An import quota restricts the amount allowed to enter the country.

        Note that export restrictions are exactly similar, only that the "quota" is then imposed by the exporting country.
    }
\end{frame}

\begin{frame}{Import Quotas}
    \showmytikz[1,2,3]{../shared/tikz/tikz-quota.tex}
    \mycomment{
        This effectively reduces the volume traded. The welfare effect can be the same as a tariff, only that the quota does not raise government revenue.
    }
\end{frame}

\begin{frame}{Import Quotas}
    \showmytikz[1,2,3,4]{../shared/tikz/tikz-quota.tex}
    \mycomment{
        However, how quotas are distributed remains in question.

        The government could use a lottery, a first-come-first-served system, or auction the quotas. If the quota do not cost the importers, they are effectively a present to them (importer surplus, $IS$).

        We see: a quota can have the same economic effect as a tariff. In the case of tariff, the price is manipulated; in the case of quota, the quantity. Since price and quantity depend on one another, the result can be the same.
    }
\end{frame}


% \section{Protectionism II: Subsidies}

\begin{frame}{Subsidies I - Subsidising Production}
    \showmytikz[1]{../shared/tikz/tikz-subsidy-trade.tex}
    \mycomment{
        Subsidies are financial assistance to producers by the government, lowering their cost of production and engaging higher output. There are two types: production subsidies and export subsidies. Here, we look at production subsidies. These lower the cost of the producer for both the world market and the internal market. The aim will be to reduce imports.
    }
\end{frame}
\begin{frame}{Subsidies I - Subsidising Production}
    \showmytikz[1,2]{../shared/tikz/tikz-subsidy-trade.tex}
    \mycomment{
    The subsidy has shifted the supply curve to the right. However, \textit{real} costs remain the same; only a part of it is now covered by the government. 

    Consumer surplus remains the same, as price and quantity are left unchanged. They will, however, now consume more local produce.

    This subsidy may reduce producer's incentives to become more cost efficient: the state shoulders the extra costs. Since the costs are not going down, producer surplus is not increased.
    }
\end{frame}

\begin{frame}{Subsidies I - Subsidising Production}
    \showmytikz[1,2,3]{../shared/tikz/tikz-subsidy-trade.tex}
    \mycomment{
        However, the producer surplus has increased due to the subsidy increasing domestic prices. 
    }
\end{frame}

\begin{frame}{Subsidies I - Subsidising Production}
    \showmytikz[1,2,3,4]{../shared/tikz/tikz-subsidy-trade.tex}
    \mycomment{
        Of course, the subsidy must have come from somewhere. If it is financed by taxes, the loss of (consumer) income will (in some markets) shift the demand curve (to the left), reducing consumer surplus. There is also an excess burden due to the costs of tax collection and subsidy payout.
        
        There is also a welfare loss: that part of the subsidy which exceeds the gain in producer surplus (dead weight loss, DWL). 
    }
\end{frame}


\begin{frame}{Subsidies II - Export Subsidy}
    \showmytikz[1]{../shared/tikz/tikz-exportsubsidy.tex}
    \mycomment{
        A country exports (already), because there are more people outside the country willing to pay the price for the goods than inside the country.    

        The government subsidises only exports: the supply curve remains the same, but for every exported unit local producers receive an extra payment.

        If the country is small, this does not affect the world price.
    }
\end{frame}

\begin{frame}{Subsidies II - Export Subsidy}
    \showmytikz[1,2]{../shared/tikz/tikz-exportsubsidy.tex}
    \mycomment{By supporting the export, producers will produce more at the higher price. The consumer inside the country will now also have to pay the higher price.}
\end{frame}

\begin{frame}{Subsidies II - Export Subsidy}
    \showmytikz[1,2,3]{../shared/tikz/tikz-exportsubsidy.tex}
    \mycomment{The size of the subsidy is larger than the `gain' of welfare, so there is a dead weight loss. 
    
    On top, of course, the government must raise the funds for the subsidy elsewhere, by taxing either another industry or through income taxes. This has of course detrimental effects on the other industries and to internal demand, which is not modelled here.

    }
\end{frame}

\begin{frame}{Subsidies II - Export Subsidy}
    \mycomment{
    So far, we have looked at a small country, so that the subsidy does not affect the world price. If, however, the quantity the country (now) exports is large, global prices might indeed fall. 
    
    As a result, consumers in the rest of the world profit from lower prices. The export subsidy then effectively "sponsors" welfare gains of consumers abroad, and results in welfare loss of producers abroad.
    }
    \showmytikz[1,2,3,4]{../shared/tikz/tikz-exportsubsidy.tex}
\end{frame}


\begin{frame}{Subsidies III - From Import to Export by Subsidy}
    \showmytikz[1]{../shared/tikz/tikz-exportsubsidy2.tex}
    \mycomment{Here, the original difference between the world price and the local cost structure is so that the country is an importer. However, the government subsidies so much that the situation  inverses - clearly a costly and inefficient affair. 
    
    }
\end{frame}

\begin{frame}{Subsidies III - From Import to Export by Subsidy}
    \showmytikz[1,2]{../shared/tikz/tikz-exportsubsidy2.tex}
    \mycomment{
    The subsidy is pushing a low domestic price so much up that the country becomes an exporter, even though it should be an importer.

    }
\end{frame}

\begin{frame}{Subsidies III - From Import to Export by Subsidy}
    \showmytikz[1,2,3]{../shared/tikz/tikz-exportsubsidy2.tex}
    \mycomment{
    A large government payout increases local producers' surplus at the costs of local consumers, but there is a large dead weight loss: the subsidy is much larger than the part which `only' shifts surplus from consumers to producers. 

    This all is under the assumption that this aggressive policy does not change the world market price.
    }
\end{frame}

\begin{frame}{Non-tariff barriers}
    \mytext{
    \begin{itemize}
        \item requirements
        \item restrictions (weapons etc.)
        \item customs delays
        \item technical barriers
        \item licensing, standards
        \item packaging, labeling, localisation
        \item sanitation rules
        \item public procurement policy
    \end{itemize}
    }
    \mycomment{Non-tariff barriers typically affect internal and external suppliers - otherwise, they would be unfair practice.    

    Among these is public procurement. According to WTO rules, and in the EU in particular, governments must procure internationally without discrimination. A noteworthy exception (or not?) is arms procurement. 
    }
\end{frame}

 